\begin{conf-abstract}
{Tracking \ce{CO2} injection to an aquifer for CCS in Iceland}
{Matthias S. Brennwald, Jonas Simon Junker, Chuan Wang, Rolf Kipfer, Anne Obermann, Martin Voigt, and Alba Zappone}
{Eawag}
{Mineral sequestration of \ce{CO2} in basalt aquifers is a promising pathway for permanent carbon storage. Within the DemoUpStorage R\&D project, \ce{CO2} was dissolved in saline groundwater and co-injected with helium as an inert tracer into a deep basalt aquifer at a pilot injection site near Helguvik, Iceland. We tested a combination of two novel geochemical and geophysical tools to monitor the subsurface \ce{CO2} dynamics. Real-time, on-site monitoring of dissolved gases was conducted over one year using a gas-equilibrium membrane-inlet mass spectrometer (GE-MIMS). Electrical resistivity tomography (ERT) was used to track changes in aquifer resistivity before and after injection. GE-MIMS measurements revealed a clear increase in He concentrations downstream of the injection point, confirming fluid transport. However, no corresponding \ce{CO2} increase was observed, indicating substantial \ce{CO2} retardation relative to the fluid transport. No anomalies in He or \ce{CO2} were detected in the overlying freshwater aquifer, indicating minimal upward migration. ERT data showed localized resistivity decreases suggesting basalt dissolution near the injection borehole. Together, GE-MIMS and ERT provided complementary, cost-effective insights into \ce{CO2} transport, reaction, and storage processes in the basalt aquifer. These tools enhance monitoring and assessment capabilities, supporting the development of safe and effective geological \ce{CO2} storage.}
\end{conf-abstract}
